\section{Summary}\label{sec:agg-creds-summary}

This chapter addresses a gap in existing anonymous credential systems.  While
later systems decentralise issuance, the issuers as a whole remain responsible
for all attributes. We instead target the setting of distinct issuers without
requiring any coordination between issuers, or indeed between the registration
authority, issuers, and verifiers at any time. The user simply interacts with
each as needed.

We achieve this through constructing an aggregate anonymous credentials system,
in which credentials issued independently by different authorities can be
combined by the user into a single compact credential and presented through one
proof. Issuers do not need to maintain state across users. We take the certified
key assumption implicit in similar constructions and make it explicit,
introducing the registration authority as its own role, and using this role to
achieve efficient multi-issuer anonymous credentials. This is the first instance
of the wider argument of this thesis. The registration authority was already
present in every scheme of this kind, as the assumption that a user's key is
certified. Naming it did not add a trusted party. It let us give that party a
single narrow capability, binding a key to a tag, and nothing more. The
authority to certify attributes stays with the issuers, who never coordinate.
The knowledge of what those attributes are stays with the user.  What the
registration authority learns is a key and a tag, and what it can do is bind one
to the other. The efficiency of the scheme comes from that division, not from
removing a participant.

We present an instantiation based on structure preserving signatures, and build
a special purpose proof for efficiency. Even with 64 attributes disclosed, both
proving and verification complete in under 10ms, with cost scaling in the number
of attributes disclosed. A concrete security analysis of the scheme would be
valuable for determining its real world security guarantees.

We briefly outline epoch-based revocation, which would address attributes with
natural expiry dates.  However, a more thorough consideration of revocation is
needed before the scheme would be able to support different attribute revocation
policies.

The anonymous credential scheme as it stands uses the issuer public keys in the
clear. As the identity of the issuers can reveal information about the user,
recent work has focused on issuer hiding. Adding the property of issuer hiding
to this protocol would be interesting future work.
